\documentclass[12pt]{article}
\usepackage[utf8]{inputenc}
\usepackage[T1]{fontenc}
\usepackage{hyperref}
\title{Kusur-Myet Language Resources for K-12 Curriculum Development}
\author{Compiled by Claude Sonnet 5 (claude-sonnet-5) for Andr\'as Kornai}
\date{}
\begin{document}
\maketitle
\section{Introduction}
Kusur-Myet (Sur--Myet; ISO 639-3 tdl) is a minor Plateau language of Bauchi and Plateau States, Nigeria, with two closely related varieties, Sur and Myet. Other names in search results include Tapshin, Nsur and Dishili. Speaker estimates in the sources found range from about 3,000--4,000 (Endangered Languages Project) to 11,000--16,000 (Joshua Project, Wikipedia). Very little material exists; this report lists only what a web search returned.
\section{Primary Educational Resources}
No primers, textbooks or school materials were found in search results.
\section{Linguistic and Reference Materials}
\begin{itemize}
\item Wikipedia overview of the Sur--Myet language: \url{https://en.wikipedia.org/wiki/Sur%E2%80%93Myet_language}
\item Decker, Danladi, Dabet, Abraham and Jonah (2021), sociolinguistic profile (Journal of Language Survey Reports), Glottolog record: \url{https://glottolog.org/resource/reference/id/656697}
\item Roger M. Blench (2023), Glottolog record: \url{https://glottolog.org/resource/reference/id/658379}
\item Endangered Languages Project entry: \url{https://endangeredlanguages.com/node/96601}
\end{itemize}
\section{Religious and Community Resources}
\begin{itemize}
\item Joshua Project language page: \url{https://joshuaproject.net/index.php/languages/tdl}
\item Joshua Project, Tapshin people group: \url{https://joshuaproject.net/people_groups/15267}
\end{itemize}
\section{Digital and Audio Resources}
\begin{itemize}
\item Global Recordings Network, Kusur-Myet: \url{https://globalrecordings.net/en/language/tdl}
\end{itemize}
\section{Curriculum Development Context}
No evidence of school use was found. Any curriculum would start from near-zero materials; the 2021 sociolinguistic profile is the best starting point for community and orthography questions.
\end{document}